\begin{statementbox}
	\textbf{\textcolor{CStatement}{条件}}

	\begin{description}[style=nextline, leftmargin=2.8em, labelsep=0.5em, itemsep=0.35em]
		\item[\textbf{\textcolor{CStatement}{条件 1（定义约定）：}}]
		      本文采用“试验次数型几何分布”：随机变量 $\xi$ 表示直到第一次成功所需的总试验次数，因此取值为 $1,2,\dots$。

		\item[\textbf{\textcolor{CStatement}{条件 2（分布定义）：}}]
		      随机变量 $\xi$ 满足
		      \[
			      P(\xi=k)=q^{k-1}p,\quad q=1-p,\; k=1,2,\dots,
		      \]
		      其中 $p\in(0,1]$ 为每次试验的成功概率，且各次试验相互独立、成功概率保持不变。
	\end{description}

	\textbf{\textcolor{CStatement}{结论}}

	\begin{description}[style=nextline, leftmargin=2.8em, labelsep=0.5em, itemsep=0.45em]
		\item[\textbf{\textcolor{CStatement}{结论一（期望公式）：}}]
		      \textbf{\textcolor{CStatement}{直观描述：}} 在上述定义下，首次成功所需总试验次数的平均值等于成功概率的倒数。
		      \[
			      E\xi = \frac{1}{p}.
		      \]

		\item[\textbf{\textcolor{CStatement}{推广结论（期望反比）：}}]
		      \textbf{\textcolor{CStatement}{直观描述：}} 成功概率 $p$ 越小，首次成功所需试验次数的平均值越大。
		      \[
			      E\xi \cdot p = 1.
		      \]
	\end{description}

	\textbf{\textcolor{CStatement}{参数限制（成功概率范围）：}} $p\in(0,1]$；当 $p=0$ 时，首次成功不可能发生，本文这种有限正整数取值的几何分布无定义；当 $p\to 0^+$ 时，$E\xi\to +\infty$。

	\medskip

	\textbf{\textcolor{CStatement}{定义提醒：}} 有些资料把几何分布定义为“首次成功前的失败次数”。若用 $X$ 表示首次成功前的失败次数，则
	\[
		P(X=k)=q^k p,\quad k=0,1,2,\dots,
	\]
	此时 $EX=\dfrac{q}{p}$，不是 $\dfrac{1}{p}$。做题时必须先确认随机变量表示的是“总试验次数”还是“失败次数”。
\end{statementbox}